\section{Introduction}
\label{sec:intro}

The cabbage stem flea beetle (CSFB, \emph{Psylliodes chrysocephala}) is a major pest of oilseed
rape (\emph{Brassica napus}). In autumn, adult beetles feed on the cotyledons and first true
leaves of young plants, leaving small round \emph{shot holes} and shallow yellow/brown
\emph{pitting}. With fewer insecticides available, breeding resistant cultivars is an important
alternative. The Res4StRes project therefore phenotypes hundreds of genotypes in field trials at
six locations. For every plot, three photographs are taken through a $0.1\,\text{m}^2$ metal
frame, and trained raters estimate the percentage of damaged leaf area by comparing the plants
with visual damage scales. Early growth stages (BBCH10--11: cotyledons, possibly the first true
leaf~\cite{bbch}) matter most, because this is when feeding damage threatens crop establishment.

Manual scoring is slow and subjective. Scoring all 8{,}946 images would take a single rater over a
year at 20 minutes per image, and when two universities (JLU and GAU) scored the same 900 images
independently, their scores correlated with only $r = 0.549$. According to the project partners,
the score folds shot holes and pitting into one number and is most likely a per-plant estimate
averaged over the plants in the image. As a result, only 2{,}747 images have scores, and only 470
images, where both raters agree within 5 percentage points, can be regarded as reliable labels.

This makes automatic scoring a \emph{weakly supervised} problem. There is one noisy scalar label
per image, no labels for individual plants or lesions, and very few reliable examples. Moreover,
a seedling covers less than 1\,\% of the image, and a shot hole only a few dozen pixels, so
simply resizing the image for a standard backbone destroys the signal. Finally, the model has to
work on new fields with different soil, lighting and plant age.

We compare three families of approaches on this problem. The first is classical computer vision:
frame detection, colour-based plant segmentation and a direct measurement of damaged pixels, plus
an RF-DETR hole/pitting detector trained on a small annotated set. The second is whole-image
regression with a frozen DINOv3 backbone, the baseline suggested in the project brief. The third
is multiple instance learning (MIL), where each image is a bag of plant regions: a frozen DINOv3
embeds each plant, a shared head scores it, and the plant scores are averaged with weights
proportional to visible plant area. We compare this with uniform pooling and attention-based MIL,
and add a pairwise ranking loss, which was our chosen research direction.

Our contributions are: (i) a leakage-safe benchmark on the 470 reliable images with plot-grouped
splits and a test set that stayed closed during development; (ii) evidence that plant-level MIL
reduces the error of the whole-image baseline by almost half, with area-weighted pooling
outperforming attention pooling; (iii) evidence that adding a ranking loss gives a small but
consistent improvement; (iv) a zero-shot evaluation on two unseen field trials and on images the
two raters disagreed on, which shows where the model fails; and (v) an analysis of which classical
components work, which do not, and why.

Section~\ref{sec:related} reviews related work, Section~\ref{sec:method} describes the methods,
Section~\ref{sec:experiments} presents the results, and Section~\ref{sec:conclusion} concludes with
limitations and recommendations for future data collection.
